\scanpage{99}
\mspage{\textcircled{2}}
On définit, si $T'$ est de $=$ relatif à $I$
\[
\begin{aligned}
\underline{\mathrm{Hom}}_\lambda(I;T,T')={}&\text{catégorie des }\inserted{\text{couples }(u,\alpha),\ \text{avec }u}\\
 &\text{morphismes } u:T_\lambda\longrightarrow T'_\lambda\ \text{\deleted{avec} isom de}\\
 &\text{commutation }\alpha
\end{aligned}\tag*{(2.4)}
\]
\diagram{\begin{tikzcd}[ampersand replacement=\&,column sep=small]
 T_\lambda \arrow[r,"u"]\arrow[d,"b_\lambda^T"']\& T'_\lambda\arrow[dl,"b_\lambda^{T'}"]\\
 \tau^I \& \deleted{\unsure{\mathrm{id}}}
\end{tikzcd}}

On a une \deleted{foncteur} 2-cat.
\[
(\lambda\text{-}I\text{-types})\lhook\joinrel\longrightarrow
\unsure{\text{(}}\underline{\mathrm{Cat}\ 2\,\mathrm{cofib}}\,(\mathrm{Cat}_{\lambda}))/\tau^{I}\tag*{(2.5)}
\]
et un 2-foncteur
\[
(\lambda\text{-}I\text{-types})\longrightarrow(\lambda\text{-types})\tag*{(2.6)}
\]
qui est $\unsure{0|1}$-fidèle.

\fbox{3} Si $I$ est réduit à un élément, $\tau^I$ est noté \deleted{$\unsure{I}$} $\tau$, donc
\[
\tau(C)=C\tag*{(3.1)}
\]
\deleted{\illegible{1} catégorie} On pose
\[
\underline{\mathrm{Hom}}_\lambda(T,\tau)\simeq R_T^\lambda\quad\deleted{\unsure{\mathrm{Rep}|(\text{ou }R_T)}}\tag*{(3,2)}
\]
et l'appelle Cat. des $\lambda$-\uline{constructions pour} $T$. On \deleted{\unsure{trouve}}\inserted{ne}
\deleted{\unsure{peut}} sait pas à priori si $R_T^\lambda$ est une $\mathcal U$-catégorie. Mais
c'est une \deleted{\unsure{cat}} $\lambda$-catégorie, et les $\lambda$-limites se calculent
``argument par argument'', donc pour toute $C\in\mathrm{Ob}\,\mathrm{Cat}_\lambda$
et $\xi\in T(C)$, le foncteur
\[
\xi_{\unsure{\xi|*}}^\lambda:R_T^\lambda\longrightarrow C\tag*{(3.3)}
\]
\deleted{commute} est un $\lambda$-foncteur.

\margin{On a (pour définir (3.3))
$(3.4)\ \left\{\begin{aligned}
 R_T^\lambda\times T(C)&\longrightarrow C\\
 \deleted{\unsure{T(C)\times R_T^\lambda\to C}}&\\
 &\text{pour }C\in\mathrm{ob}\,\mathrm{Cat}_\lambda
 \end{aligned}\right.$
i.e. $(3.5)\ R_T^\lambda\longrightarrow\underline{\mathrm{Hom}}(T(C),C)$
\deleted{Notons le foncteur \unsure{Cov} correspondant \unsure{à} \unsure{commutant} aux}
qui est un $\lambda$-foncteur.}

Bien entendu, la catégorie $R_T^\lambda$ dépend de $T$ de
façon contravariante.

\margin{Quand $I$ est donné avec $T$, on a
$(\unsure{3.6|3.5})\ I\longrightarrow\mathrm{Ob}\,R_T^\lambda$ (fonctoriel en $T\in\mathrm{Ob}(\lambda\text{-}I\text{-types})$),
qui avec (3.4) redonne (2.2), et dont la donnée \deleted{\unsure{avec}} équivaut à (2.3)
i.e. (2.2) pour tt $C$.}
% margin note written sideways
\margin{Notons aussi, de façon équivalente à (3.5)
$(3.7)\ T(C)\longrightarrow\underline{\mathrm{Hom}}(R_\lambda^T,C)$
\uline{définissant}, pour $T$ variable
$(3.8)\ T_\lambda\longrightarrow\widehat{R_\lambda^T}{}^{\lambda}=(C\longmapsto\underline{\mathrm{Hom}}(R_\lambda^T,C))$
\fbox{5} \deleted{\unsure{pour l'instant}} \unsure{nonobstant} le fait que $R_\lambda^T$ soit une $\mathcal U$-catégorie \unsure{ou}}
% slanted
\margin{\deleted{\unsure{T}} $T_\lambda$ arrive (universellement ?) dans la cat. 2-cofibrée sur $(\mathrm{Cat}_\lambda)$ \uline{représentée} par $R_T^\lambda$}

\fbox{4} \uline{Théories $\lambda$-représentables}

On dit que $T$ est $\lambda$-représentable si \deleted{la catégorie} \inserted{la catégorie}
2-cofibrée $T_\lambda$ sur $(\mathrm{Cat}_\lambda)$ est 2-repr, l'objet
$R_T$ qui la représente est appelé \uline{catégorie $\lambda$-modulaire} de $T$ : \deleted{\unsure{Elle}}
\[
T(C)\xleftarrow{\simeq}\underline{\mathrm{Hom}}_\lambda(R_T,C)\tag*{(4.1)}
\]
via une \deleted{objet} $T$-structure $\lambda$-universelle
\[
\xi_T^\lambda\in T(R_T)\tag*{(4.2)}
\]
% archive sheet number: 98
